%! Project: Design and Conduct a Survey — Unit 1, Lesson 1.8 — Guided Notes & Practice
% THE in-class document, in the MAIN IDEAS / NOTES shape (LESSON_SHAPE.md §1, ported from AP
% Statistics 2026-09-12): the vocabulary box, the hook, then ONE two-column guidednotes table.
% Each row is one idea — a label on the left; on the right one or two COMPLETE printed
% sentences, ONE large pre-drawn display the student reads or names, then a bold prompt and a
% two-across grid of numbered problems with 1.4–2.2cm of writing room. Problems run 1..19.
% THIS IS THE UNIT'S PROJECT LESSON, kept as one: the unit run forwards. Rows 1–4 rebuild the
% cycle on the Student Council's finished survey; the Guided Practice row IS the class's own
% survey — the instrument the class builds and a tally table filled live.
% DENSITY: a \blank{} only where a word or number IS the answer — the fill table of row 2
% (three broken items, the one \wordbank strip naming the bias each carries), the relative
% frequency column of row 3, problem 7 (its \wordbank), and the instrument card + tally table of the
% Guided Practice row (live data, no bank). No mid-sentence blanks; every display is
% pre-drawn; the page belongs to the student's pen.
%   I DO   : the teacher reads the definition, marks up the display, works the first problem
%            of each grid (1, 5, 8, 12).
%   WE DO  : the class works the rest of each grid, students holding the pen; the last row,
%            Guided Practice (16–19), is the class's own survey, worked entirely together.
%   YOU DO : nothing on this page — ../homework/ IS the individual practice, started in the
%            last ten minutes of the period. Its last item is the project deliverable.
% THE TRAP is problem 9 (0.36 x 900 = 324 computed correctly off a convenience sample — the
% arithmetic is right and the statistics are wrong).
% THE CRUX is problem 13 (row 4, "The Limit"): the class's 36% lands next to the council's
% 35%, and that near-agreement is luck, not evidence, because chance never chose the class —
% the hook vote is re-taken there (AFDA.DA.2j). Problem 18 is the crux transferred to the
% class's own live data.
%
% THE DATA
%   Riverbend Student Council: $3,000 grant; stratified 60 of 900 at 1/15 -> 18/16/14/12 by
%     grade (270/240/210/180); counts 21/18/15/6 -> 35%/30%/25%/10%; 0.35 x 900 = 315,
%     0.25 x 900 = 225; 42 of 120 = 35% (problem 10)
%   Last spring's class of 25 (the hook, problems 9, 13, 14): 9 club funding -> 36%;
%     0.36 x 900 = 324 — computed correctly, reported wrongly
%   Guided Practice: THIS class's live data. The key is a WORKED EXEMPLAR on a class of 25
%     answering the grant question 9/7/6/3 -> 36%/28%/24%/12% (the plan's teacher note says
%     so; the key does not).
%
% PAGE LOCKSTEP with notes_key/: the key is GENERATED from this file by templates/lesson/mkkey.py
% (spec: templates/lesson/examples/lesson08_notes_key.json) — every \blank{} becomes an \ans{},
% every \vterm{} a \vtermans{}{}, every \par\writelines{n} n \ansline{} calls, and the answer
% argument of every \pcell, \writespace and \labelbox is filled. Answer spaces are fixed
% heights, so the two files cannot drift. KEEP EVERY \pcell ON ONE SOURCE LINE — mkkey matches
% line by line.
% TABLE MECHANICS: inside a cell, `\\` ENDS THE ROW — break lines with \par. A row cannot break
% across pages: the figure gets its own sub-row (`\\` then `& ...`) and EACH GRID ROW is its own
% sub-row — close the probgrid, `\\ &`, reopen as probgrid* (no top rule).
\documentclass[12pt]{article}
\usepackage{saar-article}
\usepackage{saar-boxes}
\newcommand{\TallMath}[1]{$\displaystyle #1\rule[-1.4em]{0pt}{3.2em}$}
% Word bank strip — ONLY above a table the student fills (rows 2 and 3). Defined per document;
% the key inherits it and prints the same strip, so heights match.
\newcommand{\wordbank}[1]{\par\vspace{2pt}\noindent\fcolorbox{goldacc}{goldbg}{\parbox{\dimexpr\linewidth-2\fboxsep-2\fboxrule\relax}{\small\textbf{Word bank:} #1}}\par\vspace{4pt}}
% Vocabulary rows: a FIXED-HEIGHT row carrying the term label and OPEN writing space —
% no underline beside the term and no rule line (user correction 2026-09-06: the black
% answer line crowded the row; students write beside and below the term). The key fills
% exactly the same height, so the box cannot drift blank vs. keyed. saar-article's
% \termblank adds an inline blank and a rule line, so the pair is defined here (the key is
% generated from this file and inherits both). 2.3cm per row gives students three
% handwritten lines of room; keep key definitions to two lines.
\newcommand{\termrowheight}{1.85cm}
\newlength{\termrowinset}\setlength{\termrowinset}{7pt}
\newcommand{\vterm}[1]{%
  \par\noindent\begin{minipage}[t][\termrowheight][t]{\linewidth}%
    \vspace*{\termrowinset}%
    \textbf{\textcolor{cerulean}{#1:}}%
  \end{minipage}\par}
\newcommand{\vtermans}[2]{%
  \par\noindent\begin{minipage}[t][\termrowheight][t]{\linewidth}%
    \vspace*{\termrowinset}%
    \textbf{\textcolor{cerulean}{#1:}}\quad\ans{#2}%
  \end{minipage}\par}

\begin{document}

\pageheader{Unit 1, Lesson 1.8}{Guided Notes \& Practice}
% NAMESTRIP: no \namedateperiod here. The student writes their name once, on the
% cover this component is stapled behind.

\vspace{0.15in}

\begin{vocabbox}
\small
Fill in each term \textbf{as we name it} in the notes below.
\vterm{Survey instrument}
\vterm{Closed-response item}
\vterm{Relative frequency}
\vterm{Bar graph}
\vterm{Convenience sample}
\end{vocabbox}

\boxguard[12]
\begin{hookbox}
\small
Last spring a class of $25$ in this room asked itself the Student Council's grant question and
got $36\%$ for club funding; the council's stratified sample of $60$ had said $35\%$. One student
wrote: \emph{``We matched the council, so our class can speak for all $900$ students too.''}

\vspace{3pt}
\textbf{Is the student right?} Circle one: \quad agree \qquad disagree \quad --- and say why
you think so. We vote again at problem~13.
\par\writelines{2}
\end{hookbox}

\small
\begin{guidednotes}
% ── ROW 1: the cycle as a deliverables checklist ─────────────────────────────
\mainidea[Your project,]{The Cycle} &
A survey is an \textbf{observational study} --- you ask and you record --- so the plan is Lesson
1.1's four-stage cycle. Today it is a checklist: \textbf{each stage hands in one thing.}
\par\vspace{3pt}
\stepnum{1}\ \textbf{Formulate} --- a statistical question and its population.\quad
\stepnum{2}\ \textbf{Collect} --- a sampling plan and the finished instrument.\par
\stepnum{3}\ \textbf{Represent} --- a frequency table and a bar graph.\quad
\stepnum{4}\ \textbf{Analyze \& report} --- a report with its limits named.
\\
& {\centering
\begin{tikzpicture}[scale=0.76]
  % the council's survey, one box per stage; the student names each stage below
  \foreach \i/\txt in {
      0/{``Which of four options would the most Riverbend students choose for the \$3{,}000?''\\ Population: all 900},
      1/{A stratified sample of 60 from 900, 18/16/14/12 by grade; one item, four choices},
      2/{Counts 21 / 18 / 15 / 6 in a frequency table, then a bar graph of the percents},
      3/{``Club funding leads at 35\% --- about 315 of the 900. Chance chose the sample.''}} {
    \fill[frost, rounded corners=5pt] ({\i*3.05},0) rectangle ({\i*3.05+2.8},2.2);
    \draw[cerulean, thick, rounded corners=5pt] ({\i*3.05},0) rectangle ({\i*3.05+2.8},2.2);
    \node[font=\tiny, align=center, text width=2.0cm] at ({\i*3.05+1.4},1.1) {\txt};
  }
  \foreach \i in {0,1,2} { \draw[->, very thick, goldacc] ({\i*3.05+2.82},1.1) -- ({\i*3.05+3.03},1.1); }
  \node[font=\scriptsize\bfseries\color{cerulean}] at (5.95,2.5) {The Student Council's survey, stage by stage};
\end{tikzpicture}\par\vspace{3pt}
{\scriptsize\itshape Name the stage under each box.}\par\vspace{2pt}
\labelbox{2.2cm}{}\hspace{0.30cm}\labelbox{2.2cm}{}\hspace{0.30cm}\labelbox{2.2cm}{}\hspace{0.30cm}\labelbox{2.2cm}{}\par}
\\
& \notesprompt{Run the council's plan --- then look at your own.}
\begin{probgrid}{|Y|Y|}
\pcell{1}{Nobody in the council's survey was assigned to anything --- students were asked, and answers recorded. Observational study or experiment? How do you know?}{1.4cm}{} & \pcell{2}{Does \emph{``How should the council spend the grant?''} pass the three tests --- answers vary, a group is named, a measurement is named? Repair it if not.}{1.8cm}{} \\ \hline
\end{probgrid}
\\
& \begin{probgrid}*{|Y|Y|}
\pcell{3}{Stage 2: the council stratified by grade at $60 \div 900 = 1/15$. How many of the $270$ ninth graders and the $180$ twelfth graders does that call for?}{1.6cm}{} & \pcell{4}{Your class will ask only the students in this room. Who chose that sample --- chance, or the schedule? What is a sample like that called?}{1.6cm}{} \\ \hline
\end{probgrid}
\\ \hline
% ── ROW 2: the instrument — four rules, taught by repair ─────────────────────
\mainidea[Writing]{The Instrument} &
The \textbf{survey instrument} is the exact wording read to every person, plus its choices. A
\textbf{closed-response item} offers a fixed set of choices, so every answer lands in one
category and can be counted. \textbf{Four rules:}
\par\vspace{3pt}
\stepnum{1}\ \textbf{One idea per item.}\quad \stepnum{2}\ \textbf{Neutral wording.}\par
\stepnum{3}\ \textbf{Choices cover everyone, without overlapping} --- then the percents add to
$100\%$.\quad \stepnum{4}\ \textbf{Short.}
\\
& {\centering
\begin{tikzpicture}[scale=0.76]
  % the council's finished instrument card, annotated with the four rules
  \fill[white, rounded corners=5pt] (0,0) rectangle (8.7,3.4);
  \draw[charcoal, thick, rounded corners=5pt] (0,0) rectangle (8.7,3.4);
  \node[anchor=north west, font=\scriptsize\bfseries] at (0.25,3.22) {Student Council Grant Survey};
  \node[anchor=north west, font=\scriptsize, inner sep=1pt] at (0.25,2.72)
    {Which \textbf{one} should the grant fund? Check one.};
  \foreach \i/\txt in {0/Club funding, 1/Field trips, 2/Gym equipment, 3/Longer lunch} {
    \draw[charcoal] (0.45,{1.62-\i*0.42}) rectangle (0.70,{1.87-\i*0.42});
    \node[anchor=west, font=\scriptsize] at (0.86,{1.745-\i*0.42}) {\txt};
  }
  % callouts: the four rules, pointed at the part of the card that obeys each
  \node[font=\tiny\bfseries\color{cerulean}, anchor=west] at (9.3,3.0) {rule 1: one idea};
  \node[font=\tiny\bfseries\color{cerulean}, anchor=west] at (9.3,2.5) {rule 2: no ``don't you agree''};
  \node[font=\tiny\bfseries\color{cerulean}, anchor=west] at (9.3,2.0) {rule 4: one line long};
  \node[font=\tiny\bfseries\color{cerulean}, anchor=west, align=left, text width=2.6cm] at (9.3,1.1)
    {rule 3: four boxes --- every student fits exactly one};
  \draw[goldacc, thick, ->] (9.2,3.0) -- (8.8,2.95);
  \draw[goldacc, thick, ->] (9.2,2.5) -- (8.8,2.55);
  \draw[goldacc, thick, ->] (9.2,2.0) -- (8.8,2.35);
  \draw[goldacc, thick, ->] (9.2,1.1) -- (8.8,1.0);
\end{tikzpicture}\par}
\\
& Each item below breaks one rule. Name the rule --- and the name Lesson 1.4 gave that bias.
\wordbank{rule 1 \;\textbullet\; rule 2 \;\textbullet\; rule 3 \;\textbullet\; leading \;\textbullet\; double-barrelled \;\textbullet\; overlapping}
\footnotesize\renewcommand{\arraystretch}{1.1}
\begin{tabularx}{\linewidth}{@{} X p{1.9cm} p{3.4cm} @{}}
\toprule
\textbf{The item as written} & \textbf{Rule} & \textbf{Its 1.4 name} \\
\midrule
``Don't you agree that our underfunded clubs deserve the grant?''
  & \blank{1.5cm} & \blank{3.0cm} \\
``Should the council fund clubs and field trips?''
  & \blank{1.5cm} & \blank{3.0cm} \\
``How old are you: $14$--$16$, $16$--$18$, or $18+$?''
  & \blank{1.5cm} & \blank{3.0cm} \\
\bottomrule
\end{tabularx}\small
\\
& \notesprompt{A biased instrument is wrong in a direction --- and no arithmetic afterwards can pull it back.}
\begin{probgrid}{|Y|Y|}
\pcell{5}{Repair the leading item so it is neutral. Then say why no arithmetic done later can undo a leading item.}{1.8cm}{} & \pcell{6}{Repair the age item so a $16$-year-old has exactly one box to check. Which check on the finished percents would have caught the overlap?}{1.8cm}{} \\ \hline
\end{probgrid}
\\ \hline
% ── ROW 3: counts → percents → picture ───────────────────────────────────────
\mainidea[Counts to percents to picture]{Relative Frequency} &
A category's \textbf{relative frequency} is its count divided by the total, written as a percent;
the four add to $100\%$. ``Choice'' is \emph{categorical}, so its display is a \textbf{bar graph}:
one bar per category, bars separated, \textbf{height $=$ relative frequency}.
\par\vspace{3pt}
\textbf{7.} The council's $60$ answers came back. Complete the relative frequency column --- each count divided by $60$.
\wordbank{35\% \;\textbullet\; 30\% \;\textbullet\; 25\% \;\textbullet\; 10\% \;\textbullet\; 100\%}
\footnotesize\renewcommand{\arraystretch}{1.05}
\begin{tabularx}{\linewidth}{@{} X c p{3.2cm} @{}}
\toprule
\textbf{Choice} & \textbf{Frequency} & \textbf{Relative frequency} \\
\midrule
Club funding   & $21$ & \blank{2.6cm} \\
Field trips    & $18$ & \blank{2.6cm} \\
Gym equipment  & $15$ & \blank{2.6cm} \\
Longer lunch   & $6$  & \blank{2.6cm} \\
\midrule
\textbf{Total} & $60$ & \blank{2.6cm} \\
\bottomrule
\end{tabularx}\small
\\
& {\centering
\begin{tikzpicture}[x=0.88cm, y=0.050cm]
  % the council's bar graph, pre-drawn and pre-scaled; the student reads it
  \foreach \y in {10,20,30,40}
    \draw[linegray] (0,\y) -- (10.4,\y);
  \foreach \y in {0,10,20,30,40}
    \node[left, font=\scriptsize] at (-0.05,\y) {\y\%};
  \foreach \i/\h in {1/35, 2/30, 3/25, 4/10}
    \fill[ceruleanlight] ({2.6*\i-2.0},0) rectangle ({2.6*\i-0.6},\h);
  \foreach \i/\h in {1/35, 2/30, 3/25, 4/10}
    \node[above, font=\scriptsize\bfseries] at ({2.6*\i-1.3},\h) {\h\%};
  \foreach \i/\lab in {1/Club funding, 2/Field trips, 3/Gym equipment, 4/Longer lunch}
    \node[below, font=\scriptsize, align=center, text width=2.2cm] at ({2.6*\i-1.3},-0.5) {\lab};
  \draw[cerulean, thick] (0,42) -- (0,0) -- (10.4,0);
  \node[font=\scriptsize\bfseries\color{cerulean}, rotate=90] at (-1.3,21) {relative frequency};
\end{tikzpicture}\par\vspace{5pt}
Tallest bar: \labelbox{2.4cm}{}\quad Gaps because ``choice'' is \labelbox{2.2cm}{}\par}
\\
& \notesprompt{Analyze --- and ask who chose the sample before you multiply.}
\begin{probgrid}{|Y|Y|}
\pcell{8}{About how many of all $900$ students would choose gym equipment? Is the council allowed to say so about the school? Why?}{1.8cm}{} & \pcell{9}{Last spring's class of $25$: $9$ chose club funding. Compute the percent, then $0.36 \times 900$. \textbf{Vote first:} may the class report that second number about the school?}{1.8cm}{} \\ \hline
\end{probgrid}
\\
& \begin{probgrid}*{|Y|Y|}
\pcell{10}{Suppose the council had asked $120$ students and $42$ chose club funding. Does the club-funding bar move? What does a bar's height tell a reader --- a count, or a share?}{1.8cm}{} & \pcell{11}{\textbf{In your own words:} why must the four percents add to $100\%$ --- and which instrument rule makes that happen?}{1.8cm}{} \\ \hline
\end{probgrid}
\\ \hline
% ── ROW 4: THE CRUX — whom may the report describe? ──────────────────────────
\mainidea[Whom may you describe?]{The Limit} &
A report is four sentences: \textbf{what we asked, and of whom; how the sample was chosen, and how
many answered; what we found; what it means --- and about whom.} Sentence 4 may describe the
population only if \textbf{chance} chose the sample; a class surveying itself is a
\textbf{convenience sample}, so its sentence 4 is about the class.
\\
& {\centering
\begin{tikzpicture}[scale=0.74]
  % left: the council's stratified 60
  \fill[frost, rounded corners=6pt] (0,0) rectangle (5.9,3.0);
  \draw[cerulean, thick, rounded corners=6pt] (0,0) rectangle (5.9,3.0);
  \node[font=\scriptsize\bfseries\color{cerulean}] at (2.95,2.7) {THE COUNCIL --- $60$ of $900$};
  \node[font=\tiny, align=center, text width=4.1cm] at (2.95,2.25) {stratified by grade, \textbf{chance} chose them};
  \node[font=\normalsize\bfseries] at (2.95,1.45) {$21/60 = 35\%$};
  \node[font=\tiny\itshape] at (2.95,0.8) {may describe:};
  \node[font=\scriptsize\bfseries\color{cerulean}] at (2.95,0.38) {all $900$ students};
  % right: last spring's class of 25
  \fill[goldbg, rounded corners=6pt] (8.4,0) rectangle (14.3,3.0);
  \draw[goldacc, thick, rounded corners=6pt] (8.4,0) rectangle (14.3,3.0);
  \node[font=\scriptsize\bfseries\color{goldacc}] at (11.35,2.7) {THE CLASS --- $25$ of $900$};
  \node[font=\tiny, align=center, text width=4.1cm] at (11.35,2.25) {one room, the \textbf{schedule} chose them};
  \node[font=\normalsize\bfseries] at (11.35,1.45) {$9/25 = 36\%$};
  \node[font=\tiny\itshape] at (11.35,0.8) {may describe:};
  \node[font=\scriptsize\bfseries\color{goldacc}] at (11.35,0.38) {?};
  % the near-agreement, in the gap between the two panels
  \node[font=\scriptsize\bfseries\color{redacc}] at (7.15,1.75) {$36\% \approx 35\%$};
  \node[font=\tiny\itshape\color{redacc}] at (7.15,1.2) {luck?};
\end{tikzpicture}\par\vspace{4pt}
The class's sentence 4 may describe: \labelbox{3.2cm}{}\par}
\\
& \textbf{Near-agreement is luck, not evidence. Chance never chose the class, so it could have
missed by twenty points and nobody could tell.}
\par\vspace{2pt}
\textbf{In your own words:} why does landing next to the council's $35\%$ not earn the class the right to talk about all $900$?
\writespace{1.3cm}{}
\\
& \notesprompt{Now settle the hook vote.}
\begin{probgrid}{|Y|Y|}
\pcell{12}{Write the council's sentence 4 --- what it means for the \$$3{,}000$, and about whom.}{1.6cm}{} & \pcell{13}{\textbf{Vote again:} agree \; disagree. Does matching the council's $35\%$ make the class's $36\%$ trustworthy for all $900$? Explain.}{1.8cm}{} \\ \hline
\end{probgrid}
\\
& \begin{probgrid}*{|Y|Y|}
\pcell{14}{Write last spring's class's sentence 4. Its subject must be the people the design actually covers.}{1.6cm}{} & \pcell{15}{Name one change to the \emph{sample} that would let the class speak about all $900$ --- and one thing that change would cost.}{1.6cm}{} \\ \hline
\end{probgrid}
\\ \hline
% ── ROW 5: GUIDED PRACTICE — the class's own survey, worked live ─────────────
\mainidea[Guided practice]{Our Class Survey} &
Now we run one. The class agrees on \emph{one} question, checks it against the four rules, asks
every student in this room, and tallies. Write the item and the choices exactly as agreed.
\\
& \fcolorbox{charcoal}{white}{\parbox{\dimexpr\linewidth-2\fboxsep-2\fboxrule\relax}{%
\textbf{Our instrument} --- read word for word to every person.\par\vspace{3pt}
Item: \blank{9.2cm}\par\vspace{5pt}
(A) \blank{2.0cm}\ (B) \blank{2.0cm}\ (C) \blank{2.0cm}\ (D) \blank{2.0cm}}}
\\
& \textbf{Tally.} Ask everyone in the room; count; then divide each count by the total.
\par\vspace{2pt}
\footnotesize\renewcommand{\arraystretch}{1.05}
\begin{tabularx}{\linewidth}{@{} p{1.4cm} X p{1.6cm} p{2.0cm} @{}}
\toprule
\textbf{Choice} & \textbf{Tally} & \textbf{Count} & \textbf{Percent} \\
\midrule
(A) & & \blank{1.2cm} & \blank{1.6cm} \\
(B) & & \blank{1.2cm} & \blank{1.6cm} \\
(C) & & \blank{1.2cm} & \blank{1.6cm} \\
(D) & & \blank{1.2cm} & \blank{1.6cm} \\
\midrule
\textbf{Total} & & \blank{1.2cm} & \blank{1.6cm} \\
\bottomrule
\end{tabularx}\small
\\
& \notesprompt{We work these together.}
\begin{probgrid}{|Y|Y|}
\pcell{16}{Name the population, the sample, and $n$. Who chose it --- chance, or the schedule? What kind of sample is that?}{1.6cm}{} & \pcell{17}{Show the division for our largest category, then write our sentence 4 --- its subject is this class, not the school.}{1.6cm}{} \\ \hline
\end{probgrid}
\\
& \begin{probgrid}*{|Y|Y|}
\pcell{18}{Whose result may be reported about all $900$ --- the council's $60$, or ours? If ours lands near theirs, what does that prove?}{1.6cm}{} & \pcell{19}{Name the limit you would fix first in our design, and what fixing it would cost.}{1.6cm}{} \\ \hline
\end{probgrid}
\\ \hline
\end{guidednotes}

% ── THE NOTES END AT GUIDED PRACTICE ─────────────────────────────────────────
% There is deliberately no "you do" here: ../homework/ is the individual practice,
% started in class in the last ten minutes while the teacher circulates.

\end{document}
